% sections/counterfactual.tex — Counterfactual Simulation Engine
\section{Counterfactual Simulation}
\label{sec:counterfactual}

A growing need in audit analytics and financial risk management is the
ability to answer ``what if'' questions: \emph{What would financial
statements look like under a recession? How would a control failure
propagate through the enterprise?}  Traditional synthetic data generators
produce a single dataset from fixed parameters.  \system{} extends this
with a \emph{counterfactual simulation engine} that models causal
dependencies between financial processes, applies structured interventions,
and generates paired baseline/counterfactual datasets suitable for causal
inference research, stress testing, and audit planning.

% ─────────────────────────────────────────────────────────────────────
\subsection{Causal DAG Model}
\label{sec:cf_dag}

\system{} models the interdependencies among financial processes as a
directed acyclic graph (DAG).  Each node represents an observable or
latent parameter of the generation model, and each edge encodes the
mechanism by which a change in one parameter propagates to another.

\begin{definition}[Financial Causal DAG]
\label{def:causal_dag}
A \emph{financial causal DAG} is a tuple $G = (V, E, \tau, w)$ where:
\begin{itemize}
  \item $V = \{v_1, \ldots, v_n\}$ is a set of \emph{process nodes},
    each with a baseline value $b_i \in \R$, an optional bound interval
    $[\ell_i, u_i]$, and a category
    $c_i \in \{\text{Macro}, \text{Operational}, \text{Control},
    \text{Financial}, \text{Behavioral}, \text{Regulatory},
    \text{Outcome}\}$.
  \item $E \subseteq V \times V$ is a set of directed edges such that
    $(V, E)$ is acyclic.
  \item $\tau : E \to \mathcal{T}$ assigns a \emph{transfer function}
    to each edge, drawn from a family $\mathcal{T}$ of eight function
    types (\Cref{tab:transfer_functions}).
  \item $w : E \to [0,1]$ assigns a \emph{strength} weight to each
    edge, modulating the magnitude of the transferred effect.
\end{itemize}
Additionally, each edge carries a \emph{lag} $\lambda_e \in \N_0$
(in months), encoding the temporal delay before an upstream change
manifests downstream.
\end{definition}

The default DAG shipped with \system{} contains 17 nodes spanning
macroeconomic drivers, operational parameters, control variables, and
financial outcome metrics (\Cref{tab:dag_nodes}).

\begin{table}[t]
\centering
\caption{Causal DAG nodes with default bound intervals.  Nodes are
  organized by category; unbounded quantities use the effective range
  of the generation model.}
\label{tab:dag_nodes}
\small
\begin{tabular}{clccc}
\toprule
\textbf{ID} & \textbf{Node} & \textbf{Category}
  & \textbf{Lower} & \textbf{Upper} \\
\midrule
$v_1$    & GDP growth              & Macro       & $-0.10$ & $0.15$ \\
$v_2$    & Interest rates          & Macro       & $0.00$  & $0.20$ \\
$v_3$    & Transaction volume      & Operational & $0$     & $\infty$ \\
$v_4$    & Avg.\ transaction amount & Operational & $0$    & $\infty$ \\
$v_5$    & Control effectiveness   & Control     & $0.0$   & $1.0$ \\
$v_6$    & Fraud detection rate    & Outcome     & $0.0$   & $1.0$ \\
$v_7$    & Misstatement risk       & Outcome     & $0.0$   & $1.0$ \\
$v_8$    & Staffing level          & Operational & $0.0$   & $\infty$ \\
$v_9$    & Automation rate         & Operational & $0.0$   & $1.0$ \\
$v_{10}$ & Material cost           & Financial   & $0$     & $\infty$ \\
$v_{11}$ & Working capital         & Financial   & $-\infty$ & $\infty$ \\
$v_{12}$ & Revenue growth          & Financial   & $-1.0$  & $\infty$ \\
$v_{13}$ & Operating margin        & Financial   & $-1.0$  & $1.0$ \\
$v_{14}$ & Audit scope             & Regulatory  & $0.0$   & $1.0$ \\
$v_{15}$ & Regulatory burden       & Regulatory  & $0.0$   & $1.0$ \\
$v_{16}$ & IT investment           & Operational & $0$     & $\infty$ \\
$v_{17}$ & Process complexity      & Behavioral  & $0.0$   & $1.0$ \\
\bottomrule
\end{tabular}
\end{table}

\begin{table}[t]
\centering
\caption{Transfer function types and their mathematical forms.
  In each case, $x$ denotes the input delta from the upstream node's
  baseline.}
\label{tab:transfer_functions}
\small
\begin{tabular}{lll}
\toprule
\textbf{Type} & \textbf{Formula} & \textbf{Use Case} \\
\midrule
Linear & $ax + b$ & Proportional effects \\
Exponential & $b(1+r)^{x}$ & Compounding growth/decay \\
Logistic & $\frac{K}{1 + e^{-s(x - m)}}$
  & Saturating positive effects \\
Inverse Logistic & $\frac{K}{1 + e^{s(x - m)}}$
  & Diminishing effectiveness \\
Step & $\begin{cases} M & x > \theta \\ 0 & \text{otherwise} \end{cases}$
  & Regime switches \\
Threshold & $\min\!\big(M\frac{x-\theta}{\theta},\, S\big)$
  & Gradual activation with cap \\
Decay & $a \, e^{-rx}$ & Exponential decay \\
Piecewise & linear interpolation on $\{(x_k, y_k)\}$
  & Arbitrary empirical shapes \\
\bottomrule
\end{tabular}
\end{table}

\paragraph{Causal propagation.}
Given a set of interventions $\mathcal{I} = \{(v_j, \hat{x}_j)\}$ that
override specific node values, the engine computes the equilibrium state
of all nodes via forward propagation in topological order.
\Cref{alg:propagation} formalizes this procedure.

\begin{algorithm}[t]
\caption{Causal forward propagation with lag awareness.}
\label{alg:propagation}
\SetKwInOut{Input}{Input}
\SetKwInOut{Output}{Output}
\Input{Causal DAG $G = (V, E, \tau, w)$; interventions
  $\mathcal{I} = \{(v_j, \hat{x}_j)\}$; current month $t$}
\Output{Node values $\{x_i\}_{i=1}^{|V|}$}
\BlankLine
Compute topological order $\sigma = (\sigma_1, \ldots, \sigma_{|V|})$
  via Kahn's algorithm\;
\ForEach{$v_i \in V$}{
  $x_i \leftarrow b_i$ \tcp*{initialize to baseline}
}
\ForEach{$(v_j, \hat{x}_j) \in \mathcal{I}$}{
  $x_j \leftarrow \hat{x}_j$ \tcp*{apply direct interventions}
}
\For{$k = 1$ \KwTo $|V|$}{
  $v_i \leftarrow \sigma_k$\;
  \If{$(v_i, \cdot) \in \mathcal{I}$}{\textbf{continue}\;}
  $\Delta_i \leftarrow 0$\;
  \ForEach{edge $e = (v_p, v_i) \in E$}{
    \If{$t \geq \lambda_e$}{
      $\delta_p \leftarrow x_p - b_p$\;
      $\Delta_i \leftarrow \Delta_i + \tau_e(\delta_p) \cdot w_e$\;
    }
  }
  $x_i \leftarrow \text{clamp}\big(b_i + \Delta_i,\; \ell_i,\; u_i\big)$\;
}
\Return $\{x_i\}_{i=1}^{|V|}$\;
\end{algorithm}

The topological ordering is computed once at DAG validation time using
Kahn's algorithm~\citep{kahn1962topological} and cached.  Validation also
checks for duplicate node identifiers and dangling edge references,
returning a descriptive error if the graph is cyclic or malformed.  Node
bounds clamping (line~13 of \Cref{alg:propagation}) ensures that
propagated values remain physically plausible---for example, a fraud
detection rate cannot exceed 1.0 or drop below 0.0 regardless of
intervention severity.

% ─────────────────────────────────────────────────────────────────────
\subsection{Intervention Framework}
\label{sec:cf_interventions}

An \emph{intervention} is a structured perturbation applied to the
generation process.  \system{} defines eight intervention types,
organized by the aspect of the enterprise they target:

\begin{enumerate}
  \item \textbf{ParameterShift.} A direct change to a configuration
    value, specified via a dot-path into the YAML schema
    (e.g., \texttt{distributions.amounts.components[0].mu}).
    Supports four interpolation modes during gradual onset: linear,
    exponential, logistic, and step.

  \item \textbf{ControlFailure.} Reduces or eliminates the effectiveness
    of internal controls.  Four subtypes are modeled:
    effectiveness reduction, complete bypass, intermittent failure
    (with configurable failure probability), and delayed detection
    (with a detection lag in months).

  \item \textbf{MacroShock.} Simulates macroeconomic events such as
    recession, inflation spike, currency crisis, interest rate shock,
    commodity shock, pandemic disruption, supply chain crisis, or credit
    crunch.  Each shock type maps to pre-configured causal DAG
    node overrides; a severity multiplier scales the effect.

  \item \textbf{EntityEvent.} Targets specific entities: vendor default,
    customer churn, employee departure, new vendor onboarding,
    merger/acquisition, vendor collusion, customer consolidation, or key
    person risk.  Entity selection supports filtering by cluster,
    attribute, or explicit ID list.

  \item \textbf{ProcessChange.} Modifies business processes: approval
    threshold changes, new approval levels, system migrations, process
    automation, outsourcing transitions, policy changes, or
    reorganizations.

  \item \textbf{RegulatoryChange.} Models regulatory shifts: new standard
    adoption, materiality threshold changes, reporting requirement changes,
    or tax rate changes.

  \item \textbf{Composite.} Bundles multiple interventions into a named
    package with conflict resolution semantics (first-wins, last-wins,
    average, or error on conflict).

  \item \textbf{Custom.} Provides arbitrary config path overrides and
    downstream causal triggers for user-defined scenarios.
\end{enumerate}

\paragraph{Intervention timing.}
Each intervention specifies a timing profile consisting of a start month,
optional duration, and an \emph{onset type} that controls how the
intervention ramps in:

\begin{itemize}
  \item \textbf{Sudden:} Full effect at the start month.
  \item \textbf{Gradual:} Linear ramp from zero to full effect over a
    configurable ramp period of $r$ months.  At month $t$ within the ramp
    window:
    \begin{equation}
    \alpha(t) = \min\!\left(\frac{t - t_{\text{start}}}{r},\; 1\right)
    \label{eq:gradual_onset}
    \end{equation}
  \item \textbf{Oscillating:} Sinusoidal modulation
    $\alpha(t) = \tfrac{1}{2}(1 - \cos(\pi (t - t_{\text{start}}) / r))$,
    enabling cyclic stress patterns.
\end{itemize}

\noindent The \texttt{InterventionManager} validates timing (start month
within generation window, non-negative duration), checks intervention
bounds against node constraints, resolves config paths, and applies
priority-based conflict resolution when multiple interventions target
the same parameter.

% ─────────────────────────────────────────────────────────────────────
\subsection{Paired Generation and Diff Analysis}
\label{sec:cf_paired}

The central output of counterfactual simulation is a \emph{paired
dataset}: a baseline generated from the original configuration and a
counterfactual generated from a mutated configuration, both sharing the
same deterministic seed (\Cref{sec:arch_determinism}).  This ensures
that differences between the two datasets are attributable solely to the
intervention, not to random variation.

\paragraph{Config mutation.}
The \texttt{ConfigMutator} translates propagated causal effects into
concrete configuration changes.  It serializes the base
\texttt{GeneratorConfig} to a JSON value tree, navigates dot-paths
(including array indexing, e.g., \texttt{components[0].mu}), and
applies the new values.  After mutation, the modified configuration is
deserialized back and validated against scenario constraints.

\paragraph{Scenario engine workflow.}
The \texttt{ScenarioEngine} orchestrates the end-to-end process:
\begin{enumerate}
  \item \textbf{Load and validate} the scenario definition (interventions,
    timing, constraints).
  \item \textbf{Propagate} interventions through the causal DAG per month
    via \Cref{alg:propagation}.
  \item \textbf{Mutate} the base config using propagated node values.
  \item \textbf{Generate} baseline and counterfactual datasets from the
    original and mutated configs, respectively.
  \item \textbf{Diff} the two datasets using the \texttt{DiffEngine}.
\end{enumerate}

\paragraph{Diff analysis.}
The \texttt{DiffEngine} produces a \texttt{ScenarioDiff} comprising four
complementary views:

\begin{itemize}
  \item \textbf{ImpactSummary:} KPI-level deltas (absolute change,
    percentage change, direction) for key financial metrics, plus
    financial statement impacts (revenue, COGS, margin, net income,
    total assets, liabilities, cash flow).
  \item \textbf{RecordLevelDiff:} Per-file comparison of added, removed,
    modified, and unchanged records, with sample field-level changes.
  \item \textbf{AggregateComparison:} Statistical comparison of
    distributions, counts, and ratios across the two datasets.
  \item \textbf{InterventionTrace:} Causal attribution linking each
    observed change back to the intervention(s) that caused it, including
    the propagation path through the DAG.
\end{itemize}

Additionally, \texttt{AnomalyImpact} and \texttt{ControlImpact}
structures capture how the intervention affected anomaly injection rates
and internal control effectiveness, including whether the counterfactual
introduced new material weaknesses (per SOX~404
classification~\citep{sox2002}).

\paragraph{CLI interface.}
The scenario subsystem is exposed via four CLI subcommands:
\begin{center}
\small
\begin{tabular}{ll}
\toprule
\textbf{Command} & \textbf{Description} \\
\midrule
\texttt{datasynth-data scenario list} & List built-in and user scenarios \\
\texttt{datasynth-data scenario validate} & Validate scenario YAML \\
\texttt{datasynth-data scenario generate} & Run paired generation \\
\texttt{datasynth-data scenario diff} & Compute diff on existing outputs \\
\bottomrule
\end{tabular}
\end{center}

% ─────────────────────────────────────────────────────────────────────
\subsection{Scenario Constraints}
\label{sec:cf_constraints}

Interventions perturb generation parameters, but certain accounting
invariants must hold unconditionally---even under extreme counterfactual
scenarios.  Without enforcement, a macro shock that halves transaction
volume could inadvertently produce unbalanced journal entries or orphaned
document references, yielding a dataset that is structurally invalid
rather than merely counterfactual.

\begin{definition}[Scenario Constraint Set]
\label{def:constraints}
A \emph{scenario constraint set} $\mathcal{C}$ is a collection of
Boolean predicates over the mutated configuration $\hat{\theta}$ and
the generated dataset $\hat{D}$:
\[
\mathcal{C} = \{c_k : (\hat{\theta}, \hat{D}) \to \{\text{\textsc{true}},
  \text{\textsc{false}}\}\}_{k=1}^{|\mathcal{C}|}
\]
A counterfactual generation is \emph{valid} iff $\forall\, c_k \in
\mathcal{C}: c_k(\hat{\theta}, \hat{D}) = \text{\textsc{true}}$.
\end{definition}

\system{} enforces four built-in constraints, each enabled by default:

\begin{enumerate}
  \item \textbf{Accounting identity} (\texttt{preserve\_accounting\_identity}).
    For every journal entry $j$ in the counterfactual dataset,
    $\sum_{l \in j} l.\text{debit} = \sum_{l \in j} l.\text{credit}$
    (\Cref{eq:balanced}).

  \item \textbf{Document chain integrity}
    (\texttt{preserve\_document\_chains}).
    Every document reference (PO$\to$GR$\to$Invoice$\to$Payment) resolves
    to a valid target record, and timestamps respect causal ordering.

  \item \textbf{Period close execution}
    (\texttt{preserve\_period\_close}).
    The period-close engine (accruals, depreciation, FX translation,
    eliminations) runs to completion in every fiscal period, even if
    reduced transaction volume would otherwise leave certain accounts
    empty.

  \item \textbf{Balance coherence}
    (\texttt{preserve\_balance\_coherence}).
    The accounting equation $A = L + E$ holds at each period-end
    balance sheet.
\end{enumerate}

The \texttt{ConfigMutator} validates these constraints \emph{after}
applying mutations, rejecting any configuration that would violate them.
Users may additionally define custom constraints as range bounds on
arbitrary config paths (e.g., requiring that the anomaly injection rate
never exceeds 15\% even under stress).

This constraint mechanism ensures that counterfactual datasets remain
usable for downstream tasks---an auditor training a misstatement
detection model on the counterfactual output can trust that structural
integrity holds, and that observed differences reflect the intended
intervention rather than generation artifacts.
